\ifdefined\MOneSmall
\documentclass[acmsmall,screen]{acmart}
\else
\documentclass[manuscript,screen,review]{acmart}
\fi
\usepackage{booktabs}
\usepackage{enumitem}
\usepackage[htt]{hyphenat}
\usepackage{seqsplit}
\setcopyright{none}
\settopmatter{printacmref=false}
\acmJournal{TOMPECS}
\acmDOI{}
\acmYear{2026}
\title{Testing LLM-Serving Scheduler Ranking Portability with Trace-Driven Simulation}
\author{Soroush Vahidi}
\email{sv96@njit.edu}
\affiliation{\institution{New Jersey Institute of Technology}\city{Newark}\state{New Jersey}\country{USA}}
\begin{document}
\begin{abstract}
Comparative conclusions from simulation studies can depend on the workload,
operating condition, and metric used to produce them. We study this problem
for large language model (LLM) serving schedulers through LSSP, a
trace-driven deterministic discrete-event simulation study with a
preregistered paired selection-and-comparison design. The evaluation uses
three independently released workload sources, 120 workload windows, six
calibrated operating regions, and a fixed scheduler panel, producing
$9{,}360$ unique policy--window--region configurations and $18{,}720$
deterministic executions. Methodologically, LSSP treats ranking portability
$R(W,M,L)$ as the simulation output of interest, combines practical-margin
and confidence-interval criteria for uncertainty-aware reversal detection,
and estimates the workload sample complexity required for ranking recovery.
Cross-source rank agreement is strong between two sources (Kendall's
$\tau_b$ up to $1.0$) but lower when BurstGPT is involved (as low as
$0.55$); $3.6\%$ of pairwise comparisons ($36/990$) show statistically
supported reversals, all concentrated in one policy pair with BurstGPT on
one side. Cross-metric portability is more limited: median rank correlation
is $\tau_b=0.419$ and top-policy agreement is $31.9\%$. A selected physical
vLLM/GPU validation did not reproduce the simulator-predicted major
reversal, while the stable-control ordering qualitatively persisted.
Simulation-derived scheduler rankings should therefore be treated as
conditional on workload, operating region, metric, and model fidelity.
\end{abstract}

\ccsdesc[500]{Computing methodologies~Simulation evaluation}
\ccsdesc[300]{Computer systems organization~Cloud computing}
\keywords{LLM serving, scheduling, ranking portability, trace-driven simulation, performance evaluation}
\maketitle
\input{baseline/sections/introduction}
\input{baseline/sections/related_work}
\input{baseline/sections/study_design}
\input{baseline/sections/workloads}
\input{baseline/sections/methodology}
\input{baseline/sections/results}
\input{baseline/sections/sample_complexity}
\input{baseline/sections/real_system}
\input{baseline/sections/discussion}
\input{baseline/sections/limitations}
\input{baseline/sections/artifact}
\input{baseline/sections/conclusion}
\input{baseline/sections/declarations}
\bibliographystyle{ACM-Reference-Format}
\bibliography{references}
\end{document}
